\begin{abstract}
Mathematical reasoning remains a formidable obstacle for language models because of its inherent complexity and structured requirements. This study presents DeepSeekMath~7B, an extension of DeepSeek-Coder-Base-v1.5 7B, further pre-trained on a substantial 120B math-specific tokens derived from Common Crawl, as well as supplementary natural language and programming code. DeepSeekMath~7B attains an impressive 51.7\% accuracy on the challenging MATH competition benchmark, without employing external auxiliary tools or ensemble voting methods, thereby nearing the performance exhibited by Gemini-Ultra and GPT-4. When utilizing self-consistency with 64 generations, DeepSeekMath~7B attains 60.9\% on the MATH benchmark. The model’s prowess in mathematical reasoning is primarily a consequence of two critical design choices: (1) strategic utilization of extensive open-source web datasets via a carefully devised data filtering pipeline, and (2) the development of Group Relative Policy Optimization (GRPO)—a new Proximal Policy Optimization (PPO) variant—tailored to boost mathematical reasoning proficiency while enhancing PPO’s memory efficiency.
\end{abstract}

\section{Introduction}

Large language models (LLMs) have transformed artificial intelligence approaches to mathematical reasoning, leading to notable improvements on quantitative reasoning benchmarks \citep{MATH} as well as on geometry-related tasks \citep{trinh2024solving}. Additionally, these systems have become valuable tools for assisting individuals in the resolution of sophisticated mathematical challenges \citep{tao}. Despite these advances, the most sophisticated systems—such as GPT-4 \citep{gpt4} and Gemini-Ultra \citep{gemini}—remain inaccessible to the public, and currently available open-source alternatives fall well short of their performance.

In this work, we present DeepSeekMath, a specialized language model tailored for mathematics that demonstrates a marked advancement over existing open-source models and narrows the gap with GPT-4 in academic evaluations. Central to our approach is the construction of the DeepSeekMath Corpus, an extensive and high-quality pre-training dataset encompassing 120B tokens focused on mathematics. The corpus is sourced from the Common Crawl (CC) via a classifier built on fastText \citep{joulin2016fasttext}. For the classifier’s initial development, samples from OpenWebMath \citep{openwebmath} are utilized as positive examples, while a heterogeneous mix of unrelated web pages is selected as negative examples. Using this classifier, we extract a broader set of positive samples from CC, which are then meticulously curated with human annotation. The resulting enriched dataset is subsequently leveraged to retrain and improve the classifier’s effectiveness. Evaluation outcomes highlight the corpus's robustness, as evidenced by our base model, DeepSeekMath-Base 7B, which obtains scores of 64.2\% on GSM8K \citep{gsm8k} and 36.2\% on the MATH competition-level dataset \citep{MATH}, surpassing the performance of Minerva 540B \citep{minerva}. Moreover, given the multilingual nature of the DeepSeekMath Corpus, we observe gains on Chinese mathematical evaluations \citep{wei2023cmath,agieval}. We view our methodologies for curating mathematical datasets as a preliminary foundation for subsequent advancements, with substantial prospects for future enhancement.

DeepSeekMath-Base is constructed by initializing with DeepSeek-Coder-Base-v1.5 7B \citep{deepseek-coder}, as our experiments reveal that commencing from a code-pretrained model yields superior outcomes relative to using a general-purpose LLM. Notably, we observe that math-specific pre-training not only boosts performance on mathematical tasks but also confers broader improvements on general reasoning benchmarks such as MMLU \citep{mmlu} and BBH \citep{bbh}.

Subsequent to pre-training, we perform mathematical instruction tuning on DeepSeekMath-Base by leveraging data that reflects chain-of-thought \citep{cot}, program-of-thought \citep{pot,pal}, and tool-integrated reasoning approaches \citep{tora}. This results in DeepSeekMath-Instruct 7B, which surpasses all other models of equivalent (7B) scale and demonstrates competitive results even against open-source models with 70B parameters that have undergone instruction tuning.

In addition, we propose Group Relative Policy Optimization (GRPO), a new reinforcement learning algorithm derived from Proximal Policy Optimization (PPO) \citep{schulman2017proximal}. Distinct from standard PPO, GRPO eliminates the requirement for a critic network and instead determines the baseline using aggregate group scores, resulting in considerable savings of computational resources. By fine-tuning with a targeted subset of English instruction tuning data, GRPO markedly enhances the already-strong DeepSeekMath-Instruct model, showing significant performance gains across both in-domain (GSM8K: 82.9\% $\rightarrow$ 88.2\%, MATH: 46.8\% $\rightarrow$ 51.7\%) and out-of-domain math benchmarks (e.g., CMATH: 84.6\% $\rightarrow$ 88.8\%) throughout reinforcement learning. Furthermore, we propose a unified framework to interpret various methods—such as Rejection Sampling Fine-Tuning (RFT) \citep{yuan2023scaling}, Direct Preference Optimization (DPO) \citep{dpo}, PPO, and GRPO—revealing that these approaches can be viewed as variations or simplifications of RL strategies. We carry out comprehensive empirical studies—contrasting online and offline training, outcome versus process supervision, as well as single-turn and iterative RL—aimed at unraveling the core mechanisms underpinning this framework. Finally, we elucidate the factors driving RL-induced improvements in instruction-tuned models, and suggest avenues for advancing RL effectiveness based on our unified perspective.

\subsection{Contributions}

This work presents advancements in large-scale mathematical pre-training and investigates reinforcement learning methods through systematic exploration and analysis.

\textbf{Math Pre-Training at Scale}

\begin{itemize}[topsep=0pt]
    \item We demonstrate that the openly available Common Crawl resource encompasses significant mathematical content. Through a carefully engineered data selection procedure, we curate the DeepSeekMath Corpus—a dataset containing 120B tokens specifically filtered for mathematics, making it approximately seven times larger than the Minerva \citep{minerva} dataset of mathematical web pages and about nine times the size of OpenWebMath \citep{openwebmath}.

    \item Our DeepSeekMath-Base 7B pre-trained model matches the performance of Minerva 540B \citep{minerva} on relevant benchmarks, providing evidence that scaling up parameter count is not the sole determinant of mathematical reasoning proficiency; a more compact architecture, when pre-trained on higher-quality data, can also exhibit competitive results.

    \item We present observations from a suite of mathematics-focused training experiments. In particular, pre-training on code prior to math-specific data enhances model performance in solving mathematical tasks, both when using and not using external tools. This partially addresses a persistent research query: \textit{Does code training bolster reasoning skills?} Our experiments suggest a positive effect, especially regarding mathematical reasoning.

    \item Contrary to widespread practice, supplementing pre-training with arXiv papers—often the standard in mathematics-oriented studies—yields negligible gains across all evaluated mathematical benchmarks in this work.
\end{itemize}

\textbf{Exploration and Analysis of Reinforcement Learning}

\begin{itemize}[topsep=0pt]
    \item We present Group Relative Policy Optimization (GRPO), a reinforcement learning framework that is both resource-efficient and high-performing. In contrast to methods like Proximal Policy Optimization (PPO), GRPO omits the need for a critic model by estimating the baseline through aggregated group scores, which leads to a substantial reduction in computational demands during training.
    \item Our results indicate that applying GRPO substantially improves the capabilities of the instruction-tuned model DeepSeekMath-Instruct, relying exclusively on instruction-tuning datasets. Additionally, we note gains in out-of-domain generalization throughout the reinforcement learning phase.
    \item We introduce a cohesive conceptual framework that integrates techniques such as RFT, DPO, PPO, and GRPO. Through thorough empirical analysis—including comparisons of online and offline training, outcome versus process supervision, and single-turn versus iterative reinforcement learning—we systematically examine the foundational components of this approach.
    \item Leveraging this unified framework, we analyze the underlying mechanisms that contribute to reinforcement learning's efficacy, and delineate several avenues for advancing more robust reinforcement learning strategies for LLMs.
\end{itemize}

\subsection{Summary of Evaluations and Metrics}

\begin{itemize}[topsep=0pt]
    \item \textbf{English and Chinese Mathematical Reasoning}: We perform exhaustive evaluations on our models using diverse English and Chinese mathematical benchmarks that span a spectrum from elementary to university-level problems. The English datasets comprise GSM8K \citep{gsm8k}, MATH \citep{MATH}, SAT \citep{llemma}, OCW Courses \citep{minerva}, and MMLU-STEM \citep{mmlu}, while the Chinese datasets feature MGSM-zh \citep{mgsm}, CMATH \citep{wei2023cmath}, Gaokao-MathCloze \citep{agieval}, and Gaokao-MathQA \citep{agieval}. Evaluation criteria include the models’ abilities to produce independent, tool-free solutions in text form, as well as their proficiency in solving problems using Python.

    On English evaluation sets, DeepSeekMath-Base achieves performance on par with the closed-source Minerva 540B \citep{minerva}, and consistently outperforms all available open-source base models (including Mistral 7B \citep{mistral} and Llemma-34B \citep{llemma}), regardless of their exposure to math pre-training, frequently by a notable margin. Especially in the Chinese evaluation benchmarks, DeepSeekMath-Base exhibits pronounced superiority, attributable in part to our strategy of incorporating high-quality multilingual math data, rather than adhering to the English-exclusive pre-training paradigm seen in prior works \citep{minerva,llemma}. Enhanced by mathematical instruction tuning and reinforcement learning, DeepSeekMath-Instruct and DeepSeekMath-RL set new standards, becoming the first open-source models to achieve greater than 50\% accuracy on the rigorous, competition-level MATH dataset.
\end{itemize}

\item \textbf{Formal Mathematics}: The performance of DeepSeekMath-Base is assessed on the informal-to-formal theorem proving task as described in \citep{dsp_proof}, utilizing miniF2F \citep{minif2f} with Isabelle \citep{isabelle} as the designated proof assistant. In these few-shot autoformalization tasks, DeepSeekMath-Base achieves notable results.
\item \textbf{Natural Language Understanding, Reasoning, and Code}: For a well-rounded evaluation of the model's capabilities in general comprehension, logical reasoning, and programming, DeepSeekMath-Base is tested on several benchmarks. The Massive Multitask Language Understanding (MMLU) suite \citep{mmlu} measures accuracy across 57 multiple-choice topics, while BIG-Bench Hard (BBH) \citep{bbh} features 23 difficult tasks primarily involving multi-step reasoning. Coding skills are assessed with HumanEval \citep{codex} and MBPP \citep{mbpp}, both commonly employed in code model evaluation. Results indicate that math-centric pre-training enhances both linguistic reasoning and understanding.

\section{Math Pre-Training}

\subsection{Data Collection and Decontamination} This section details the construction of the DeepSeekMath Corpus sourced from Common Crawl. As illustrated in Figure \ref{fig:data_collect}, we describe an iterative methodology for amassing a large-scale mathematical dataset from Common Crawl, beginning with a limited yet high-quality math-related seed dataset. Notably, this systematic approach can be extended to other fields, such as computer programming.

The process begins by selecting OpenWebMath \citep{openwebmath}—a vetted set of high-caliber mathematical texts from the web—as the seed dataset. Using this collection, we develop a fastText classifier \citep{joulin2016fasttext} to identify additional mathematical webpages resembling those in OpenWebMath. For training, we sample 500,000 instances from the seed corpus as positive examples and 500,000 webpages from Common Crawl as negative examples. Training employs an open-source tool\footnote{\url{https://fasttext.cc}}, setting the embedding size at 256, a learning rate of 0.1, maximum n-gram length to 3, a minimum occurrence count of 3, and the number of epochs to 3. To manage the enormous scale of Common Crawl, URL-based deduplication and near-duplicate filtering reduce the pool to 40B HTML webpages. With the trained fastText model, we extract mathematics-focused pages from this deduplicated set. To ensure content quality, retrieved webpages are ranked using fastText’s predicted scores, with only the top-performing subset retained. The optimal dataset size is selected through pre-training studies across various token counts: 40B, 80B, 120B, and 160B. For the first iteration, the decision is made to keep the top 40B tokens.

Following the initial data collection phase, a significant number of mathematical web pages are left uncollected. This is primarily due to the limited diversity among the positive training examples used by the fastText model. To address this issue, we expand the seed dataset by sourcing further mathematical websites, thereby enabling us to better optimize the fastText model. Our approach begins by partitioning the entire Common Crawl dataset according to disjoint domains, where each domain encompasses web pages that share the same base URL. For each domain, we determine the proportion of web pages captured in the initial collection round. If a domain has more than 10\% of its pages included in the initial harvest, we mark it as math-related (for instance, \url{mathoverflow.net}).

Next, within these math-related domains, we manually label URLs corresponding to mathematical content (e.g., \url{mathoverflow.net/questions}). Web pages under these URLs that were not collected previously are then incorporated into the seed set. Through this iterative enrichment, we supply the fastText model with a broader array of positive examples, enhancing its capability to retrieve mathematical data in future iterations. After conducting four rounds of collection, the resulting dataset comprises 35.5M mathematical web pages and 120B tokens. During the fourth collection cycle, we observe that approximately 98\% of all data had already been gathered by the end of the third iteration, prompting us to conclude the collection process.

To mitigate the risk of benchmark data contamination, we adopt the strategy outlined in \cite{deepseek-coder}: filtering out web pages that feature questions or answers derived from major English and Chinese mathematical benchmarks such as GSM8K \citep{gsm8k}, MATH \citep{MATH}, CMATH \citep{wei2023cmath}, and AGIEval \citep{agieval}. The filtering procedure stipulates that any passage containing a 10-gram substring matching exactly any portion of these evaluation benchmarks is excluded from the mathematical training data. For benchmark items shorter than 10 grams but consisting of at least 3 grams, we use precise matching to identify and remove contaminated web pages.

\subsection{Validating the Quality of the DeepSeekMath~Corpus}

To assess the quality of the DeepSeekMath~Corpus in comparison to other contemporary math-centric datasets, we conduct a series of pre-training experiments. The reference math corpora include:

\begin{itemize}[topsep=0pt]
    \item \textbf{MathPile} \citep{mathpile}: a collection aggregating 8.9B tokens from various sources, such as textbooks, Wikipedia, ProofWiki, CommonCrawl, StackExchange, and arXiv—with over 85\% of content derived from arXiv;
    \item \textbf{OpenWebMath} \citep{openwebmath}: a corpus of 13.6B tokens sourced from CommonCrawl and curated for mathematical material;
    \item \textbf{Proof-Pile-2} \citep{llemma}: an aggregate comprising OpenWebMath, 10.3B tokens of mathematical code from AlgebraicStack, and 28.0B tokens from arXiv papers. In line with the methodology in \cite{llemma}, we employ an arXiv:Web:Code ratio of 2:4:1 during experiments on Proof-Pile-2.
\end{itemize}

\subsubsection{Training Setting}
\label{sec:quality-policy}
We perform mathematics-specific training on a broadly pre-trained language model consisting of 1.3B parameters, which follows the architecture of the DeepSeek LLMs \citep{deepseek-llm}, referred to here as DeepSeek-LLM 1.3B. Each mathematical dataset is used independently for 150B-token training runs. All training is executed utilizing the resource-efficient HAI-LLM framework \citep{haillm}. Aligning with DeepSeek LLM training procedures, optimization employs AdamW \citep{adamW} with hyperparameters $\beta_1=0.9$, $\beta_2=0.95$, and $\mathrm{weight\_decay}=0.1$. The learning rate follows a multi-phase schedule: starting with a warmup over the first 2,000 steps to peak, decreasing to 31.6\% of the maximum at 80\% completion of training, and further down to 10.0\% by 90\%. The peak learning rate is set at 5.3e-4. Training uses a batch size of 4M tokens and a context window of 4K tokens.

\subsubsection{Evaluation Results}

\textbf{The DeepSeekMath~Corpus distinguishes itself by its superior quality, multilingual coverage, and substantial scale.}

\begin{itemize}[topsep=0pt]
    \item \textbf{High-quality}: 
    The effectiveness of models trained on different mathematical corpora is assessed via performance on 8 downstream mathematical tasks using few-shot chain-of-thought prompting \cite{cot}. Table \ref{tab:corpora_comparison} demonstrates that the DeepSeekMath~Corpus consistently leads to stronger results. Additionally, as shown in Figure \ref{fig:corpora_comparisons}, models trained with DeepSeekMath outperform those trained on Proof-Pile-2 after 50B tokens (corresponding to a single epoch of Proof-Pile-2), underscoring the higher average data quality in DeepSeekMath.
    \item \textbf{Multilingual}:
    Data in the DeepSeekMath~Corpus spans multiple languages, with English and Chinese constituting the majority. As indicated in Table \ref{tab:corpora_comparison}, pre-training on DeepSeekMath~Corpus boosts performance in mathematical reasoning tasks for both English and Chinese. This contrasts with current mathematical datasets, which mainly focus on English and show limited—and sometimes detrimental—impact on Chinese mathematical reasoning.
    \item \textbf{Large-scale}:
    The DeepSeekMath~Corpus exceeds other mathematical datasets by several multiples in size. Figure \ref{fig:corpora_comparisons} reveals that DeepSeek-LLM 1.3B exhibits a steeper and more sustained improvement trajectory when trained on DeepSeekMath~Corpus. By comparison, other baseline corpora are significantly smaller, are exposed to the model for multiple epochs during training, and result in rapid performance stagnation.
\end{itemize}

\subsection{Training and Evaluating DeepSeekMath-Base 7B}

This section presents DeepSeekMath-Base 7B, a foundational model designed with advanced reasoning skills, particularly tailored for mathematical tasks. The model initialization employs DeepSeek-Coder-Base-v1.5 7B \citep{deepseek-coder} as its starting point and is further trained on a corpus comprising 500B tokens. The dataset allocation includes 56\% from the DeepSeekMath Corpus, 4\% from AlgebraicStack, 10\% sourced from arXiv, 20\% from code on Github, with the final 10\% consisting of natural language samples from Common Crawl in both English and Chinese. Training is performed primarily following the procedure described in Section \ref{sec:quality-policy}, with modifications: the learning rate's maximum is set to 4.2e-4 and the batch size to 10M tokens.

To thoroughly evaluate DeepSeekMath-Base 7B, we examine its proficiency in three main mathematical tasks: generating independent, comprehensive solutions; addressing problems utilizing auxiliary tools; and engaging in formal theorem proving. In addition to its mathematical reasoning, we also report on the model’s broader capabilities, encompassing natural language comprehension, logical inference, and coding aptitude.

\paragraph{Mathematical Problem Solving with Step-by-Step Reasoning}

For assessing DeepSeekMath-Base's aptitude in mathematical problem solving, we employ few-shot chain-of-thought prompting strategies \citep{cot} over eight different benchmarks available in both English and Chinese. These benchmarks test a range of mathematical reasoning skills, including quantitative reasoning (for instance, GSM8K \citep{gsm8k}, MATH \citep{MATH}, and CMATH \citep{wei2023cmath}) as well as multiple-choice assessments (such as MMLU-STEM \citep{mmlu} and Gaokao-MathQA \citep{agieval}), spanning subject matter from basic arithmetic to university-level mathematics.

Table \ref{tab:base_math_cot} illustrates that DeepSeekMath-Base 7B consistently achieves top performance across all eight benchmarks compared to other open-source base models, such as the popular generalist Mistral 7B \citep{mistral} and the newly introduced Llemma 34B \citep{llemma}, which received specialized mathematical training on Proof-Pile-2 \citep{llemma}. Remarkably, DeepSeekMath-Base surpasses prior open-source base models by more than 10\% absolute accuracy on the MATH dataset, which features competition-level problems, and it also exceeds the results of Minerva 540B \citep{minerva}—a much larger, closed-source model derived from PaLM \citep{palm} and further trained on mathematical corpora.

\paragraph{Mathematical Problem Solving with Tool Use} We assess the program-assisted mathematical reasoning abilities of models on GSM8K and MATH through few-shot program-of-thought prompting \citep{pot,pal}. In this setup, models are instructed to address problems by generating Python code, utilizing libraries such as \textit{math} and \textit{sympy} for handling complex computations. The final program output serves as the model’s answer. Table \ref{tab:base_math_pot_proof} demonstrates that DeepSeekMath-Base 7B establishes a new state-of-the-art among open-source models, outperforming Llemma 34B.

\paragraph{Formal Mathematics}

Automating formal proof construction enhances both the precision and efficiency of mathematical reasoning, an area that has garnered increasing interest. We test DeepSeekMath-Base 7B on informal-to-formal proof generation as defined in \citep{dsp_proof}, which entails producing a formalized proof from an informal description, its formal counterpart, and an informal justification. Experiments are conducted on the miniF2F benchmark \citep{minif2f}, covering Olympiad-level formal mathematics, where we generate Isabelle formal proofs using few-shot prompts. Consistent with \cite{dsp_proof}, our approach produces proof outlines that are then expanded by applying the automated prover Sledgehammer \citep{sledgehammer}. Table \ref{tab:base_math_pot_proof} shows that DeepSeekMath-Base 7B achieves robust results in the task of proof autoformalization.

\paragraph{Natural Language Understanding, Reasoning, and Code}

We benchmark the model’s proficiency in natural language understanding with MMLU \citep{mmlu}, its reasoning skills with BBH \citep{bbh}, and its coding performance on HumanEval \citep{codex} and MBPP \citep{mbpp}. According to Table \ref{tab:understanding_reasoning_code}, DeepSeekMath-Base 7B substantially improves over its predecessor, DeepSeek-Coder-Base-v1.5 \citep{deepseek-coder}, on MMLU and BBH, indicating that enhanced math-focused training benefits both linguistic comprehension and reasoning abilities. Further, by incorporating code data for continuous training, DeepSeekMath-Base 7B sustains the coding capability demonstrated by DeepSeek-Coder-Base-v1.5 on both coding evaluations. Overall, DeepSeekMath-Base 7B significantly outperforms the generalist Mistral 7B \citep{mistral} across all three tasks involving reasoning and code.

\section{Supervised Fine-Tuning}

\subsection{SFT Data Curation}

We compile a comprehensive instruction-tuning dataset for mathematics, encompassing both English and Chinese problems across multiple domains and complexity levels. Each problem is paired with an accompanying solution in one of several formats: chain-of-thought (CoT) \citep{cot}, program-of-thought (PoT) \citep{pot,pal}, or tool-integrated reasoning as outlined in \citep{tora}. Our final training set comprises 776K examples.

\begin{itemize}[topsep=0pt]
    \item \textbf{English mathematical datasets}: Tool-integrated annotations are provided for GSM8K and MATH problems, supplemented by a curated portion of MathInstruct \citep{MathInstruct} and the Lila-OOD training split \citep{lila}, in which problems have solutions represented in CoT or PoT styles. The English corpus incorporates a variety of mathematical areas, including algebra, probability, number theory, calculus, and geometry.
    \item \textbf{Chinese mathematical datasets}: We gather a range of K-12 mathematical problems in Chinese that span 76 distinct subfields, such as linear equations, annotating the solutions using both CoT and tool-integrated methodologies.
\end{itemize}

\subsection{Training and Evaluating DeepSeekMath-Instruct 7B}

This section details the instruction tuning procedure of DeepSeekMath-Instruct 7B, which builds on DeepSeekMath-Base. For training, examples are concatenated at random up to a maximum context window of 4K tokens. The model is optimized for 500 steps, employing a batch size of 256 and a fixed learning rate of 5e-5.

To measure the mathematical capabilities of the model, we conduct evaluations with and without tool assistance, using four benchmarks for quantitative reasoning in both English and Chinese. Performance comparisons are made with leading contemporaneous models:

\begin{itemize}[topsep=0pt]
    \item \textbf{Closed-source models}: This group features (1) the GPT series, with GPT-4 \citep{gpt4} and GPT-4 Code Interpreter~\footnote{\url{https://openai.com/blog/chatgpt-plugins\#\#code-interpreter}} noted for their advanced capabilities, (2) Gemini Ultra and Pro \citep{gemini}, (3) Inflection-2 \citep{inflection-2}, (4) Grok-1~\footnote{\url{https://x.ai/model-card}}, and recent releases by Chinese technology companies, including (5) Baichuan-3~\footnote{\url{https://www.baichuan-ai.com}}, and (6) the latest GLM-4~\footnote{\url{https://open.bigmodel.cn/dev/api\#glm-4}} belonging to the GLM series \citep{glm}.
\end{itemize}

These models serve general-purpose applications, and the majority have undergone various alignment procedures.

\item \textbf{Open-source models} comprise: foundational models such as (1) DeepSeek-LLM-Chat 67B \citep{deepseek-llm}, (2) Qwen 72B \citep{qwen}, (3) SeaLLM-v2 7B \citep{seallm}, and (4) ChatGLM3 6B \citep{chatglm3}, alongside models tailored for mathematical proficiency including (5) InternLM2-Math 20B~\footnote{\url{https://github.com/InternLM/InternLM-Math}}, which extends InternLM2 through additional mathematical pretraining and subsequent instruction tuning, (6) Math-Shepherd-Mistral 7B, which utilizes PPO training \citep{schulman2017proximal} on Mistral 7B \citep{mistral} via a process-supervised reward framework, (7) WizardMath models \citep{wizardmath}, which augment mathematical reasoning capabilities in Mistral 7B and Llama-2 70B \citep{llama2} by applying evolve-instruct methods (i.e., instruction tuning with AI-generated prompts) and PPO training, primarily utilizing data from GSM8K and MATH, (8) MetaMath 70B \citep{metamath}, built by fine-tuning Llama-2 70B on an enriched version of GSM8K and MATH, (9) ToRA 34B \cite{tora}, obtained by fine-tuning CodeLlama 34B to perform tool-augmented mathematical reasoning, and (10) MAmmoTH 70B \citep{MathInstruct}, which involves instruction tuning Llama-2 70B on MathInstruct.

Referencing Table \ref{tab:sft_rl_math}, in the evaluation scenario prohibiting tool usage, DeepSeekMath-Instruct 7B achieves impressive results for stepwise reasoning. In particular, on the advanced MATH benchmark, this model outperforms all open-source alternatives and most closed-source systems (e.g., Inflection-2 and Gemini Pro) by no less than 9\% absolute margin, including models that are both significantly larger in parameter count (such as Qwen 72B) and those specifically improved through math-centric reinforcement learning (such as WizardMath-v1.1 7B). Although DeepSeekMath-Instruct attains competitive accuracy relative to the Chinese proprietary models GLM-4 and Baichuan-3 on MATH, it still trails behind GPT-4 and Gemini Ultra.

When the evaluation protocol allows for integration of natural language reasoning and programming-based tool use, DeepSeekMath-Instruct 7B reaches nearly 60\% accuracy on MATH, establishing a new high among open-source models. On other relevant benchmarks, this model’s performance rivals that of DeepSeek-LLM-Chat 67B, which was the previous leading open-source model despite its tenfold parameter advantage.

\section{Reinforcement Learning}

\subsection{Group Relative Policy Optimization}
Reinforcement learning (RL) has demonstrated substantial improvements in LLMs’ mathematical reasoning capabilities after initial Supervised Fine-Tuning (SFT) \citep{wang2023math,wizardmath}. Here, we detail our proposed approach: Group Relative Policy Optimization (GRPO), designed to be both efficient and effective.

\subsubsection{From PPO to GRPO}
Proximal Policy Optimization (PPO) \citep{schulman2017proximal} is a widely adopted actor-critic algorithm in RL-based fine-tuning of LLMs \citep{ouyang2022training}. Its optimization objective involves maximizing the following surrogate function:

\begin{equation}
\footnotesize
    \mathcal{J}_{PPO}(\theta) = \mathbb{E}{[q \sim P(Q), o \sim \pi_{\theta_{old}}(O|q)]} \frac{1}{|o|} \sum_{t=1}^{|o|} \min \left[ \frac{\pi_\theta(o_{t} | q, o_{<t})}{\pi_{\theta_{old}}(o_{t} | q, o_{<t})} A_{t}, \text{clip} \left( \frac{\pi_\theta(o_{t} | q, o_{<t})}{\pi_{\theta_{old}}(o_{t} | q, o_{<t})}, 1 - \varepsilon, 1 + \varepsilon \right)  A_{t} \right] ,
\end{equation}

Here, $\pi_{\theta}$ denotes the current policy model, while $\pi_{\theta_{old}}$ refers to the previous policy. The variables $q$ and $o$ represent questions and model outputs, respectively, both drawn from the question set and the old policy $\pi_{\theta_{old}}$. The parameter $\varepsilon$ is a hyper-parameter related to clipping, introduced by PPO to stabilize training dynamics. The advantage term, $A_t$, is calculated using Generalized Advantage Estimation (GAE) \citep{gae}, which utilizes rewards $\{r_{\ge t}\}$ and an estimated value function $V_{\psi}$. As a result, PPO necessitates training a separate value function model alongside the main policy, and to prevent the reward model from being excessively optimized, it is standard practice to introduce a per-token KL divergence penalty—relative to a reference model—into the reward at every generation step \citep{ouyang2022training}:

\begin{equation}
    r_{t} = r_\varphi(q, o_{\le t}) - \beta \log\frac{\pi_{\theta}(o_{t}|q, o_{<t})}{\pi_{ref}(o_{t}|q, o_{<t})},
\label{eq:PPO-reward}
\end{equation}

where $r_\varphi$ denotes the reward model, $\pi_{ref}$ is the reference (typically the initial SFT model), and $\beta$ specifies the KL penalty's strength.

Since the value function in PPO is generally parameterized by a neural network of similar scale to the policy model itself, this approach leads to significant increases in memory and computational demands. Further, during reinforcement learning updates, this value function serves as a baseline to reduce estimator variance. Notably, in large language models, the reward model often only evaluates the reward at the final output token, which poses difficulties for accurately training a value function across all positions. To remedy this, we introduce Group Relative Policy Optimization (GRPO), depicted in Figure \ref{fig:grpo}. GRPO forgoes the separate value function approximator required by PPO, opting instead to use the mean reward from multiple responses to the same prompt as a baseline. Specifically, for each prompt $q$, GRPO samples a set of $G$ outputs $\{o_1, o_2, \ldots, o_G\}$ using $\pi_{\theta_{old}}$, then maximizes the following objective:

\begin{equation}
\footnotesize
\begin{split}
    \mathcal{J}_{GRPO}(\theta) &= \mathbb{E}{[q \sim P(Q), \{o_i\}_{i=1}^G \sim \pi_{\theta_{old}}(O|q)]}  \\
    & \frac{1}{G}\sum_{i=1}^G\frac{1}{|o_i|} \sum_{t=1}^{|o_i|} \left\{ \min \left[ \frac{\pi_\theta(o_{i,t} | q, o_{i,<t})}{\pi_{\theta_{old}}(o_{i,t} | q, o_{i,<t})} \hat{A}_{i,t}, \text{clip} \left( \frac{\pi_\theta(o_{i,t} | q, o_{i,<t})}{\pi_{\theta_{old}}(o_{i,t} | q, o_{i,<t})}, 1 - \varepsilon, 1 + \varepsilon \right)  \hat{A}_{i,t} \right] - \beta \mathbb{D}_{KL}\left[\pi_{\theta} || \pi_{ref}\right]\right\} ,
\end{split}
\label{eq:GRPO-obj}
\end{equation}

In this formulation, $\varepsilon$ and $\beta$ act as hyper-parameters, and the advantage $\hat{A}_{i,t}$ is computed from the rewards associated only with outputs within the same sample group; further elaboration on this process is provided in the subsequent sections. The group-relative advantage calculation adopted by GRPO matches the structure of many reward modeling procedures, since reward models are usually built upon datasets of direct comparisons among answers to a single query. Additionally, GRPO incorporates the KL divergence regularization as a direct term in the loss, rather than embedding it within the reward calculation as in PPO, thereby simplifying the computation of $\hat{A}_{i,t}$. Distinct from the token-level KL penalty used in (\ref{eq:PPO-reward}), we estimate the KL divergence by applying the unbiased estimator of \citep{kl_approx}:

\begin{equ

\begin{algorithm}[t]
  \small
  \caption{Iterative Group Relative Policy Optimization}
  \textbf{Input} initial policy model $\pi_{\theta_{\text{init}}}$; reward models $r_\varphi$; task prompts $\mathcal{D}$; 
  hyperparameters $\varepsilon$, $\beta$, $\mu$
  \begin{algorithmic}[1]
    \State Initialize the policy model: $\pi_\theta \leftarrow \pi_{\theta_{\text{init}}}$
    \For{iteration = 1, \dots, I}
       \State Set the reference model $\pi_{ref} \leftarrow \pi_{\theta}$
      \For{step = 1, \dots, M}
      \State Draw a mini-batch $\mathcal{D}_b$ from $\mathcal{D}$
      \State Update the previous policy model: $\pi_{\theta_{old}} \leftarrow \pi_{\theta}$ 
      \State For every question $q$ in $\mathcal{D}_b$, sample $G$ responses $\{o_i\}_{i=1}^G \sim \pi_{\theta_{old}} (\cdot \mid q) $
      \State Evaluate each sampled output $o_i$ with the reward model $r_{\varphi}$, obtaining rewards $\{r_i\}_{i=1}^{G}$
      \State For each $o_i$, estimate the group relative advantage $\hat{A}_{i,t}$ for the $t$-th token
      \For{GRPO iteration = 1, \dots, $\mu$}
        \State Adjust policy model $\pi_{\theta}$ to maximize the GRPO objective as specified by Equation \ref{eq:GC-GRPO}
      \EndFor
    \EndFor \State Continuously update $r_\varphi$ with replay-based training. 
    \EndFor 
  \end{algorithmic}
  \textbf{Output} $\pi_\theta$
  \label{alg:iter-grpo}
\end{algorithm}

\subsubsection{Outcome Supervision RL with GRPO} More formally, for each question $q$, a set of candidate outputs $\{o_1, o_2, \cdots, o_G\}$ are generated using the preceding policy model $\pi_{\theta_{old}}$. The outputs are then scored using a reward model, generating $G$ reward values $\mathbf{r}=\{r_1, r_2, \cdots, r_G\}$. These reward values are standardized by subtracting the batch mean and dividing by the batch standard deviation. With outcome supervision, the normalized reward is assigned at the sequence’s end, such that all tokens in an output receive the same normalized reward as their advantage, that is, $\hat{A}_{i, t} = \widetilde{r}_i = \frac{r_i- {\rm mean}(\mathbf{r})}{{\rm std}(\mathbf{r})}$. Policy parameters are then updated by maximizing the objective outlined in equation (\ref{eq:GRPO-obj}).

\subsubsection{Process Supervision RL with GRPO} The outcome supervision approach assigns rewards solely at output completion, which may be inadequate for guiding models through intricate mathematical reasoning. Inspired by \cite{wang2023math}, we also consider process supervision, which attaches rewards to the completion of individual reasoning steps. Specifically, for a given prompt $q$ and sampled responses $\{o_1, o_2, \cdots, o_G\}$, a process-based reward model scores each step in every output, leading to reward sets $\mathbf{R} = \{ \{r_1^{index(1)},\cdots,r_1^{index(K_1)}\}, \cdots,  \{r_G^{index(1)},\cdots,r_G^{index(K_G)}\} \}$, with $index(j)$ denoting the terminal token of step $j$, and $K_i$ the total steps for $o_i$. These step-wise rewards are standardized analogously: $\widetilde{r}_i^{index(j)} = \frac{r_i^{index(j)} - {\rm mean(\mathbf{R})}}{{\rm std(\mathbf{R})}}$.
Process supervision then determines each token’s advantage as the cumulative sum of future step-wise normalized rewards, i.e., $\hat{A}_{i, t} = \sum_{index(j) \ge t} \widetilde{r}_i^{index(j)}$,
and model parameters are optimized via maximization of the objective in equation (\ref{eq:GRPO-obj}).

\subsubsection{Iterative RL with GRPO} As reinforcement learning advances, the previously trained reward model may become inadequate for effectively guiding the evolving policy model. To address this, we experiment with iterative RL utilizing GRPO. Algorithm \ref{alg:iter-grpo} outlines our procedure, where, in each iteration of GRPO, we regenerate training datasets for the reward model using fresh samples produced by the policy model. The reward model is then incrementally updated with a replay strategy, ensuring 10\% of prior data is retained in each training batch. Subsequently, we designate the current policy model as the new reference model, proceeding to further train the policy model using the newly updated reward model.

\subsection{Training and Evaluating DeepSeekMath-RL}

Our reinforcement learning experiments are conducted on DeepSeekMath-Instruct 7B. For training, we use questions in a chain-of-thought format, sourced from GSM8K and MATH datasets within the SFT data, resulting in a dataset of approximately 144K examples. We omit all other SFT questions to better study how RL affects evaluation benchmarks for which the training data is absent during the RL phase. The reward model training set is built following the method described in \citep{wang2023math}. The reward model is initialized with DeepSeekMath-Base 7B and trained with a learning rate of 2e-5. In GRPO, the policy model uses a learning rate of 1e-6, and a KL penalty coefficient of 0.04. Each question is associated with $64$ sampled responses. The maximum sequence length and the training batch size are both set to 1024. After each stage of exploration, the policy model undergoes a single update. The evaluation of DeepSeekMath-RL 7B uses the same benchmarks as DeepSeekMath-Instruct 7B. For this setup, tasks like GSM8K and MATH that employ chain-of-thought reasoning are categorized as in-domain, while all other evaluation sets are treated as out-of-domain tasks.

Table \ref{tab:sft_rl_math} presents a comparison between open- and closed-source models that employ both chain-of-thought and tool-based reasoning, evaluated on English and Chinese datasets. Our observations are as follows: 1) When applying chain-of-thought reasoning, DeepSeekMath-RL 7B achieves accuracy rates of 88.2\% on GSM8K and 51.7\% on MATH. This places it ahead of all open-source models ranging from 7B to 70B parameters, as well as surpassing most closed-source competitors. 2) Importantly, DeepSeekMath-RL 7B’s training exclusively involves chain-of-thought instruction tuning data from GSM8K and MATH, and it is initialized from DeepSeekMath-Instruct 7B. Despite having access to a limited training corpus, it consistently exceeds DeepSeekMath-Instruct 7B’s performance across every metric, indicating the significant impact of reinforcement learning.

\section{Discussion}
This section outlines key insights gained from our experiments involving both pre-training and reinforcement learning.

\subsection{Lessons Learnt in Pre-Training}

We begin by discussing lessons from the pre-training stage. Unless specified otherwise, we employ the experimental configurations described in Section \ref{sec:quality-policy}. Additionally, when referencing the DeepSeekMath Corpus herein, it pertains to the 89B-token collection produced during the second round of data acquisition.

\subsubsection{Code Training Benefits Mathematical Reasoning}

The notion that code-focused training may bolster a model’s reasoning capabilities is widespread yet lacks thorough empirical support. Here, we provide initial evidence within the scope of mathematical reasoning: code training enhances model performance on mathematical tasks, both in the presence and absence of tool assistance.

To evaluate the effect of code training on mathematical reasoning abilities, we compared the outcomes under two training regimens: a two-stage training process and a single-stage training approach.

\textbf{Two-Stage Training}

\begin{itemize}[topsep=0pt]
    \item \textbf{Code Training for 400B Tokens $\rightarrow$ Math Training for 150B Tokens}: DeepSeek-LLM 1.3B undergoes initial training on 400B code tokens, followed by further training with 150B math tokens;
    \item \textbf{General Training for 400B Tokens $\rightarrow$ Math Training for 150B Tokens}: For comparison, the model is also pre-trained on 400B general-domain tokens (sourced from DeepSeek-AI’s large-scale general dataset) before 150B math tokens, enabling a direct analysis of the comparative value that code tokens provide versus general tokens for fostering mathematical reasoning abilities.
\end{itemize}

\textbf{One-Stage Training}

\begin{itemize}[topsep=0pt]
    \item \textbf{Math Training for 150B Tokens}: DeepSeek-LLM 1.3B is trained from scratch solely using 150B math tokens;
    \item \textbf{Training on a mixture of 400B Code Tokens and 150B Math Tokens}: Recognizing that sequential code-then-math training negatively affects coding skills, we also examine the effects of simultaneously training on a dataset combining 400B code tokens with 150B math tokens, to evaluate whether this joint approach can still elevate mathematical reasoning while lessening catastrophic forgetting.
\end{itemize}

\paragraph{Results}
Tables \ref{tab:code-math} and \ref{tab:code-math-general-eval} present the downstream results associated with each training configuration.

Integrating code training confers clear advantages for program-aided mathematical reasoning, as observed in both two-stage and one-stage paradigms. Specifically, Table \ref{tab:code-math} indicates that pretraining exclusively with code already substantially elevates model competence in resolving GSM8K and MATH challenges via Python, and that subsequent math-focused training drives further gains. Notably, in the one-stage scheme, blending code and math tokens diminishes the extent of catastrophic forgetting typically encountered in the two-stage process, while also augmenting coding proficiency (see Table \ref{tab:code-math-general-eval}) and supporting program-driven mathematical reasoning (Table \ref{tab:code-math}).

Furthermore, code-based pretraining enhances the model's intrinsic mathematical reasoning capabilities—independent of tool usage. During two-stage training, code exposure initially imparts moderate performance boosts and enhances the efficacy of subsequent math-centric training, ultimately yielding optimal results. In contrast, conducting one-stage training with both code and math tokens mixed appears to hinder tool-free mathematical reasoning, potentially due to the limited capacity of DeepSeek-LLM 1.3B to assimilate both types of knowledge effectively within a single-stage learning process.

\subsubsection{Limited Effect of ArXiv Papers on Mathematical Reasoning}

Although arXiv papers are frequently used as part of pre-training data for mathematics models \citep{minerva,gpt-f,llemma,mathpile}, a thorough investigation into their efficacy in enhancing mathematical reasoning has not yet been carried out. Surprisingly, our empirical findings indicate that including arXiv papers does little to advance mathematical reasoning abilities. To assess this, we trained models of different capacities—specifically, DeepSeek-LLM 1.3B and DeepSeek-Coder-Base-v1.5 7B \citep{deepseek-coder}—using two distinct arXiv-based datasets:

\begin{itemize}[topsep=0pt]
    \item \textbf{MathPile} \citep{mathpile}: an 8.9 billion-token dataset constructed with heuristic-based data cleaning and filtering, with more than 85\% comprised of scientific arXiv articles;
    \item \textbf{ArXiv-RedPajama} \citep{redpajama}: a 28.0 billion-token collection of arXiv LaTeX source files processed to remove preambles, comments, macros, and bibliography entries.
\end{itemize}

We trained DeepSeek-LLM 1.3B on 150B tokens and DeepSeek-Coder-Base-v1.5 7B on 40B tokens from each corpus. Consistently, models trained exclusively on arXiv papers failed to show measurable gains—and sometimes even exhibited diminished performance—across a variety of mathematical tasks of varying complexity. The evaluation spanned datasets focusing on quantitative problem solving such as GSM8K and MATH (see Table \ref{tab:arxiv-cot}), multiple-choice STEM problems such as those in MMLU-STEM (also Table \ref{tab:arxiv-cot}), as well as formal theorem-proving tasks in miniF2F (Table \ref{tab:arxiv_atp}).

Nonetheless, these findings are not definitive and come with several caveats. Notably, we have yet to examine:

\begin{itemize}[topsep=0pt]
    \item The effect of arXiv-based data on more specialized mathematical tasks not covered in our evaluations, such as the transformation of formal mathematics into informal language (the informalization of theorems and proofs);
    \item The influence of combining arXiv material with additional types of training data;
    \item Whether larger models could unlock benefits from arXiv content not visible at smaller scales.
\end{itemize}

Therefore, further studies are warranted to comprehensively understand the contributions of arXiv papers in mathematical model pre-training.

\subsection{Insights into Reinforcement Learning}
\subsubsection{Toward a Unified Framework} Here, we propose a general framework to interpret several training approaches—including SFT, RFT, DPO, PPO, and GRPO—and present experimental results exploring how different aspects of this framework influence outcomes. Typically, for such training algorithms, the gradient with respect to model parameters $\theta$ is formalized as follows:

\begin{equation}
    \nabla_{\theta}\mathcal{J}_{\textcolor{red}{\mathcal{A}}}(\theta) = \mathbb{E}[\underbrace{(q,o) \sim \textcolor{red}{\mathcal{D}}}_{Data \ Source}]\left( \frac{1}{|o|} \sum_{t=1}^{|o|}  \underbrace{GC_{{\mathcal{A}}}(q, o, t, \textcolor{red}{\pi_{{rf}}})}_{Gradient \ Coefficient}  \nabla_{\theta}\log \pi_{\theta}(o_t | q, o_{<t})\right).
\label{eq:objective}
\end{equation}

The formulation comprises three main elements: (1) \textit{Data Source $\mathcal{D}$}, specifying the corpus utilized for model training; (2) \textit{Reward Function $\pi_{{rf}}$}, which provides the evaluative feedback essential for learning; and (3) \textit{Algorithm $\mathcal{A}$}, responsible for integrating the dataset and reward to yield the gradient coefficient $GC$, dictating the strength and direction of parameter updates (penalty or reinforcement). Within this generalized framework, we examine a selection of paradigmatic approaches:

\begin{itemize}
    \item  \textbf{Supervised Fine-tuning (SFT)}: This approach involves adapting a pretrained model using manually curated SFT datasets via supervised learning.
    \item \textbf{Rejection Sampling Fine-tuning (RFT)}: RFT further adapts the SFT model by training on outputs generated by the SFT model itself, selecting only those completions that satisfy correctness criteria for each prompt.
    \item \textbf{Direct Preference Optimization (DPO)}: DPO enhances the SFT model by fine-tuning on a set of augmented outputs, optimizing using a pairwise preference-driven loss.
    \item \textbf{Online Rejection Sampling Fine-tuning (Online RFT)}: Unlike standard RFT, Online RFT commences from an SFT-initialized policy and continually improves it with outputs sampled and filtered from the evolving online policy during training.
    \item \textbf{PPO/GRPO}: Here, the policy model is initialized with the SFT weights and is then optimized using reinforcement learning on samples drawn from the current, evolving policy model.
\end{itemize}

Table \ref{tab:data_policy} presents a summary of the constituent elements of the aforementioned methods. For an expanded exposition of the derivation procedures, consult Appendix \ref{app:analysis-rl}.

\paragraph{Observation about Data Source} The categorization of data sources is twofold: online sampling and offline sampling. Online sampling involves collecting training samples directly from the ongoing outputs generated by the live policy model during training. Conversely, offline sampling sources its training examples from the outputs of the pre-trained SFT model. In this framework, RFT and DPO operate using offline data collection, whereas Online RFT and GRPO utilize data collected via the online paradigm.

Referencing Figure \ref{fig:rl-analysis},
it is evident that Online RFT achieves notably better results than RFT across two evaluation benchmarks. Early in the training process, performance between Online RFT and RFT is comparable; however, as training proceeds, Online RFT distinctly surpasses RFT, underscoring the benefits of adopting online sampling strategies. This phenomenon is reasonable: in the initial training phases, the actor and SFT models yield highly similar outputs, leading to data samples with minimal variation. As training advances, the divergence between the actor’s samples and those from the SFT model becomes more pronounced, and leveraging freshly sampled online data delivers more pronounced improvements.

\paragraph{Observation about Gradient Coefficient} When updating model parameters, each algorithm transforms the reward signal into a gradient coefficient. In this context, reward functions are distinguished as either ‘Rule’-based or ‘Model’-based. A Rule-based reward assesses response quality according to the factual correctness of an answer, while a Model-based reward employs a learned reward model to evaluate and score responses. Notably, the training set for the reward model is curated through rule-based judgments. Equations \ref{eq:GC-RFT} and \ref{eq:GC-GRPO} illustrate a salient distinction between GRPO and Online RFT: GRPO modulates its gradient coefficient as a function of the reward outputted by the reward model, permitting finer control over the reinforcement and penalization of different responses according to their scores. Online RFT, on the other hand, applies a uniform reinforcement to all correct responses and does not explicitly penalize incorrect ones, lacking the differentiated gradient adjustment found in GRPO.

As illustrated in Figure \ref{fig:rl-analysis}, GRPO achieves superior results compared to online RFT, underscoring the effectiveness of adjusting both positive and negative gradient coefficients. Moreover, GRPO+PS outperforms GRPO+OS, which highlights the advantages of utilizing fine-grained, step-aware gradient coefficients. Our investigation also covers iterative RL, where we perform two iterations in our experiments. Figure \ref{fig:iter-rl} reveals that iterative RL considerably enhances performance, particularly during the first iteration.

\subsubsection{Why RL Works?} This study applies reinforcement learning to a subset of instruction tuning data, leading to significant performance gains over the base instruction tuning model. To shed light on the mechanisms underlying these improvements, we compare the Pass@K and Maj@K accuracies for both Instruct and RL models across two benchmarks. According to Figure \ref{fig:maj-pass}, while RL provides a boost to Maj@K accuracy, Pass@K remains largely unaffected. This suggests that the benefit of RL stems from reinforcing correct answers among the TopK responses, rather than fundamentally advancing core capabilities; \textbf{that is, the improvement appears to be due to increased prominence of correct outputs within TopK rather than foundational skill enhancement.} Likewise, \citep{wang2023making} identified a \textbf{misalignment problem} affecting reasoning tasks in SFT models, and have shown that their performance can be improved by various preference alignment approaches \citep{yuan2023rrhf,song2023preference,wang2023making}.

\subsubsection{How to Achieve More Effective RL?}
\label{sec:effec_rl}
Our results indicate that RL is highly effective in mathematical reasoning scenarios. Additionally, we introduce a unified framework for interpreting a variety of prominent training approaches; within this framework, all methods are seen as forms of either direct or approximate RL techniques. As detailed in Equation \ref{eq:objective}, there are three primary components: Data Source, Algorithm, and Reward Function. We outline several promising directions for future research concerning these elements.

\paragraph{Data Source} The data source serves as the foundational input for all training strategies. Specifically in RL contexts, we define the data source as consisting of unlabeled questions alongside outputs sampled from the policy model. In our current work, only questions derived from the instruction tuning phase, paired with basic nucleus sampling for output generation, are utilized. We hypothesize that this choice of data source may partly explain why our RL pipeline primarily enhances Maj@K performance. Moving forward, we aim to investigate RL pipelines using out-of-distribution prompts and to incorporate \textbf{advanced sampling (decoding) strategies}, including those inspired by tree-search algorithms \citep{yao2023tree}. Furthermore, the adoption of \textbf{efficient inference techniques} \citep{xia-etal-2023-speculative,leviathan2023fast,kwon2023efficient,xia2024unlocking}, which are pivotal for effective exploration within policy models, remains a crucial area for future work.

\paragraph{Algorithms} Algorithms utilize both data and reward feedback to compute gradient coefficients for model parameter updates. As reflected in Equation \ref{eq:objective}, prevailing approaches fundamentally \textbf{TRUST} the guidance of the reward function when adjusting the conditional probability assigned to specific tokens. Nonetheless, it is infeasible to guarantee the consistent accuracy of reward signals, particularly in highly intricate tasks. For instance, the PRM800K datasets \citep{lightman2023let}, even though meticulously labeled by skilled annotators, are reported to contain nearly 20\% erroneous annotations\footnote{\url{https://github.com/openai/prm800k/issues/12\#issuecomment-1728491852}}. In light of these challenges, this paper examines reinforcement learning algorithms that demonstrate robustness in the face of unreliable reward signals. We argue that \textbf{WEAK-TO-STRONG} \citep{burns2023weak} alignment strategies have the potential to reshape learning algorithm foundations.

\paragraph{Reward Function} The reward function delivers essential feedback for training. In the context of RL, it is commonly instantiated as a neural reward model. We identify three pivotal aspects for advancing reward models: 1) \textbf{Strengthening the generalization capacity of the reward model.} A robust reward model should successfully extrapolate to out-of-distribution tasks and novel decoding outputs; otherwise, RL might only reinforce existing LLM distributions without genuine performance gains; 2) \textbf{Capturing and expressing model uncertainty.} Quantifying uncertainty could play a critical mediating role between less-reliable reward models and weak-to-strong learning frameworks; 3) \textbf{Developing efficient, high-fidelity process reward models} capable of yielding detailed, granular feedback throughout the reasoning trajectory \citep{lightman2023let,wang2023math}.

\section{Conclusion, Limitation, and Future Work}

In this work, we introduce DeepSeekMath, which surpasses all other open-source systems on the competitive MATH benchmark and rivals several closed-source counterparts in performance. The model builds upon DeepSeek-Coder-v1.5 7B, undergoing continued training on 500B tokens, with a notable inclusion of 120B mathematics-focused tokens extracted from Common Crawl. Through a comprehensive ablation analysis, we determine that web-based resources present significant value as sources of high-quality mathematical datasets, whereas arXiv data does not contribute as much as anticipated. We also propose Group Relative Policy Optimization (GRPO), an alternative to Proximal Policy Optimization (PPO), demonstrating that GRPO substantially boosts mathematical reasoning abilities while reducing memory requirements. Experimental findings reveal that GRPO remains beneficial even after DeepSeekMath-Instruct 7B attains strong performance on relevant benchmarks. Furthermore, we offer a unified conceptual framework to interpret various methods and highlight promising directions to enhance reinforcement learning approaches.

Despite these advancements, DeepSeekMath~exhibits certain shortcomings: its abilities in areas like geometry and theorem proving remain inferior compared to closed models. For example, our preliminary tests reveal the model's inability to resolve triangle and ellipse-related tasks, possibly reflecting biases introduced during the pre-training or fine-tuning data selection. Additionally, due to its model size, DeepSeekMath~lags behind GPT-4 in few-shot learning; while GPT-4's performance improves with few-shot examples, DeepSeekMath~displays similar results for both zero-shot and few-shot tasks. Looking ahead, we intend to refine our data selection mechanisms to curate a higher-quality pre-training corpus. Additionally, we will further investigate avenues outlined in Section \ref{sec:effec_rl} to advance reinforcement learning methodologies for large language models.

\end{document}